% !Mode:: "TeX:UTF-8"
\chapter{附录A：实验问卷与访谈提纲}

\section{问卷说明}
本问卷用于评估两种导览条件在知识讲述支持、情境响应能力与注意力引导效果三个维度上的表现。实验采用 A/B 两组交叉验证方式进行，受试者在每轮体验结束后分别填写一次问卷，并根据当次体验情况对每个条目进行 1--5 分评分（1 表示“非常不同意”，5 表示“非常同意”）。

\section{知识讲述支持维度条目}
\begin{enumerate}
    \item 系统能够围绕同一展品持续展开讲解。
    \item 系统的回答能够与我的提问保持连贯，而不是零散的信息。
    \item 系统提供的信息有助于我形成对展品内容的整体理解。
\end{enumerate}

\section{情境响应能力维度条目}
\begin{enumerate}
    \item 系统的回答能够贴合我当前关注的展品。
    \item 在连续对话中，系统能够记住前面的交流内容。
    \item 系统的讲解能够随着参观进程的变化进行调整。
\end{enumerate}

\section{注意力引导效果维度条目}
\begin{enumerate}
    \item 我能够感受到数字人在主动引导我关注展品。
    \item 数字人的动作、视线或语气有助于我判断当前讲解重点。
    \item 数字人的表达更像导览引导，而不仅是信息播报。
\end{enumerate}

\section{访谈提纲}
\begin{enumerate}
    \item 你觉得系统在展品知识讲述方面最有帮助的地方是什么？
    \item 在连续导览过程中，你是否感到系统能够理解你当前关注的对象和上下文？
    \item 数字人的语音、动作与视线是否帮助你更容易把握讲解重点？请举例说明。
    \item 两种导览方式相比，你更倾向于哪一种？主要原因是什么？
    \item 你认为系统最需要优先改进的部分是什么？
\end{enumerate}
